\chapter{Emergency Diagnostic Imaging Pathways}
\label{app:imaging_algorithms}

This appendix summarizes high-yield imaging pathways for time-critical emergencies. Protocols must be adapted to patient stability and institutional resources.

\begin{longtable}{@{}p{0.22\textwidth}p{0.24\textwidth}p{0.46\textwidth}@{}}
\caption{Rapid imaging pathways for selected medical and surgical emergencies.}
\label{tab:app_emergency_pathways} \\
\toprule
\textbf{Emergency} & \textbf{First-Line Imaging} & \textbf{Key Decision Steps} \\
\midrule
\endfirsthead
\toprule
\textbf{Emergency} & \textbf{First-Line Imaging} & \textbf{Key Decision Steps} \\
\midrule
\endhead
Acute ischemic stroke & NCCT head $\rightarrow$ CTA head/neck ($\pm$ CTP) & Exclude hemorrhage on NCCT immediately for thrombolysis eligibility; assess ASPECTS and LVO on CTA; use CTP or DWI--FLAIR MRI for extended/unknown-window selection \\
Suspected pulmonary embolism & CTPA (or V/Q / leg Doppler if CTA contraindicated) & Assess clinical pretest probability (Wells/PERC) and D-dimer first in low/intermediate risk; evaluate RV/LV ratio on CTPA for right heart strain \\
Acute aortic syndrome & ECG-gated CTA chest/abdomen/pelvis & Non-contrast phase first (to detect intramural hematoma), then arterial-phase CTA to classify Stanford Type A vs.\ Type B \\
Polytrauma (unstable) & eFAST + portable CXR/pelvic X-ray & Do not delay resuscitation or laparotomy for CT in hemodynamically unstable patients; stable patients proceed to pan-CT \\
Acute abdominal pain (perforation/obstruction) & CT abdomen/pelvis (or upright CXR + US if CT unavailable) & Look for extraluminal gas, transition point, bowel-wall ischemia, or free fluid \\
Suspected cord compression / cauda equina & Urgent whole-spine MRI & Evaluate canal stenosis, epidural mass/hematoma/abscess, and cord signal change; urgent surgical/radiation consultation \\
\bottomrule
\end{longtable}
